% ---------------------------------------------------------------------------
% Hauptteil, Teil 2: Theorie (Theoriedarstellung, Ausgangsannahmen, Hypothesen)
% ---------------------------------------------------------------------------
\chapter{Theoretischer Rahmen}
\label{cha:theorie}

Der theoretische Teil der Arbeit erläutert zunächst, anhand welcher Theorien in dem Artikel von Kanas und Kosyakova die Auswirkungen der lokalen Sprachkurse auf die Arbeitsmarktintegration beschrieben werden. Es wird erläutert, wie die Hypothesen des replizierten Artikels anhand der theoretischen Ansätze abgeleitet und wie versucht wurde, sie zu erklären. Anschließend wird im Abschnitt „Theoretische Weiterentwicklung“ erläutert, inwiefern dieser reproduzierte Artikel theoretisch erweitert bzw. ergänzt wird.

In dem reproduzierten Artikel wird anhand von vier verschiedenen theoretischen Ansätzen erläutert, warum ein höheres lokales Kursangebot die Integration von Geflüchteten in den Arbeitsmarkt beeinflusst. Die erste dieser Theorien ist die von Becker entwickelte Humankapitaltheorie. Die Humankapitaltheorie geht davon aus, dass Sprachkurse dank besserer Sprachkenntnisse und spezifischer Informationen über den Arbeitsmarkt des Aufnahmelandes die Produktivität steigern. Die Signaling-Theorie hingegen ist eine weitere Theorie, die erklärt, dass die durch Sprachkurse erworbenen Zertifikate bei fehlenden Informationen die Lernfähigkeit einer Person belegen. Die Grundannahme von Collins’ Social-Closure- bzw. Credentialing-Theorie ist, dass Sprachzertifikate eine Schlüsselrolle dabei spielen, Barrieren beim Zugang zum Arbeitsmarkt zu beseitigen. 

Abschließend wird die Theorie des Sozialkapitals – eine Theorie, die versucht, die Integration in den Arbeitsmarkt anhand von aufgebauten Netzwerken und sozialen Beziehungen zu erklären – zur Erläuterung der Forschungsfrage herangezogen. Als theoretische Weiterentwicklung wird die Komplementarität des Humankapitals nach Chiswick als Erweiterung des theoretischen Konzepts erläutert.

\section{Humankapitaltheorie}
\label{cha:human}

Die Humankapitaltheorie ist eine ökonomische Theorie, die auf der Grundidee basiert, dass das Wissen und die Fähigkeiten von Menschen in der Zukunft einen ökonomischen Ertrag erbringen können \parencite{becker1964human}. Der grundlegende Mechanismus dieser Theorie erläutert, dass Humankapital die Produktivität steigern und somit dem Individuum einen ökonomischen Ertrag bringen kann.

Das Wissen und die Fähigkeiten der Menschen werden aus ökonomischer Sicht betrachtet. Der Begriff „Humankapital“ beinhaltet das von Menschen erworbene Wissen und die Fähigkeiten. Becker erläutert, dass Humankapital ebenso wie physisches Kapital die Produktivität des Individuums steigert. Demnach wird Humankapital als das Wissen und die Fähigkeiten definiert, die die Produktivität der Menschen steigern \parencite[7--10]{becker1964human}.

Der Begriff der Investition spielt bei der Entstehung von Humankapital eine wichtige Rolle. Eine Investition bezeichnet Tätigkeiten, die Menschen gegenwärtig ausüben und die in der Zukunft einen wirtschaftlichen Nutzen bringen können. Gemäß der Humankapitaltheorie ist eine Investition mit Kosten verbunden. Die Kosten der Investition stehen in einem Zusammenhang mit dem künftigen Ertrag. Das Individuum trägt die Kosten in der Gegenwart, jedoch wird davon ausgegangen, dass diese Investition dem Individuum in der Zukunft einen ökonomischen Nutzen bringen kann. Insgesamt steigern Investitionen in das Humankapital das Humankapital des Individuums, indem sie dessen Wissen und Fähigkeiten fördern. Es wird angenommen, dass mit steigendem Humankapital die Produktivität des Individuums zunehmen kann und dies in der Zukunft zu höheren ökonomischen Erträgen führen kann \parencite[7--10]{becker1964human}.

Gemäß Beckers Theorie des Humankapitals kann der Spracherwerb als eine Form von Humankapital verstanden werden. Die Teilnahme an Sprachkursen lässt sich als eine Investition der Geflüchteten in ihr Humankapital verstehen. Es ist anzunehmen, dass Geflüchtete durch die Teilnahme an Sprachkursen in Zukunft über bessere Deutschkenntnisse verfügen werden. Die Teilnahme der Geflüchteten an lokalen Sprachkursen in Deutschland kann mit Kosten wie Zeit, Geld und Aufwand verbunden sein. Durch die Teilnahme an Sprachkursen erworbene Deutschkenntnisse können die Produktivität von Geflüchteten steigern und ihre Beschäftigung auf dem Arbeitsmarkt fördern \parencite{becker1964human}.

Eine ähnliche Verbindung zwischen Sprachkursen, Humankapital und Arbeitsmarktintegration wird auch von \textcite{Kanas:2023} hergestellt. Das lokale Angebot an Sprachkursen kann dazu beitragen, dass Migrant:innen ihre aufnahmelandspezifischen Kenntnisse und Fähigkeiten weiterentwickeln können. Durch die Teilnahme an Sprachkursen können Kenntnisse und Fähigkeiten erworben werden, die die Produktivität auf dem Arbeitsmarkt steigern können, beispielsweise bei Vorstellungsgesprächen oder bei der Kommunikation mit Arbeitskolleg:innen, wodurch Geflüchtete im Beschäftigungskontext von einem Nutzen profitieren könnten.

\section{Signaling-Theorie}
\label{cha:Signal}

Ein weiterer theoretischer Ansatz, der die Auswirkungen der Förderung lokaler Deutschkurse auf die Beschäftigung zu erklären versucht, ist Signaling-Theorie~\parencite{Spence1973}. Spences Signaling-Theorie ist ein theoretisches Modell, das das Problem erklärt, wie Arbeitgebende die Produktivität und das Leistungspotenzial von Arbeitssuchenden einschätzen, wenn sie darüber nur unvollständige Informationen besitzen \parencite[55--56]{Spence1973}.

Arbeitgebende hingegen können die Produktivität und das Potenzial des Arbeitssuchenden vor der Einstellung nicht vollständig einschätzen. Aus diesem Grund basiert dieser theoretische Ansatz auf der Informationsasymmetrie zwischen Arbeitgebenden und Arbeitssuchenden. In Spences Modell verfügen Arbeitgebende nur über unvollständige Informationen über die Produktivitätsfähigkeiten des Arbeitssuchenden. Für Arbeitgebende stehen daher beobachtbare Merkmale im Vordergrund. Laut Spence lassen sich beobachtbare Merkmale aus zwei verschiedenen Perspektiven erklären. Sie werden in zwei Kategorien unterteilt: Indizes und Signale. Indizes bezeichnen Merkmale, die der Person nicht zu kontrollieren und nicht zu verändern sind, während Signale als Merkmale definiert werden, die der Person zu kontrollieren und zu verändern stehen. Diese Unterscheidung bildet die Grundlage für Spences Definition des Begriffs „Signal“ \parencite[357--358]{Spence1973}.

Der Erwerb eines Sprachzertifikats wird hingegen als potenzielles Signal angesehen. Eine Person, die ein Sprachzertifikat erwirbt, kann Arbeitgebenden durch dieses Zertifikat die Möglichkeit geben, Rückschlüsse auf ihre Trainingsfähigkeit und Motivation zu ziehen. Dem Modell zufolge dient der Erwerb eines Sprachzertifikats Arbeitgebenden als Hinweis auf die Fähigkeiten oder das Potenzial der Person. Auf diese Weise können Arbeitgebende Rückschlüsse auf die Motivation, Produktivität und das Potenzial der Person ziehen \parencite[357--359]{Spence1973}.

Nach Spences Modell ist das Signalisieren mit Kosten verbunden. Diese Kosten müssen nicht unbedingt monetärer Natur sein. Es handelt sich auch um nicht-monetäre Kosten wie Aufwand, Mühe und Zeit. Während der Erwerb eines Zertifikats für produktive Personen mit geringeren Kosten verbunden ist, sind für Personen mit geringer Produktivität höhere Kosten erforderlich. Auf diese Weise können Arbeitgebende Rückschlüsse auf die nicht direkt beobachtbaren Eigenschaften der Person ziehen \parencite[359]{Spence1973}.

Auch wenn die durch Sprachkurse erworbenen Zertifikate die Leistungsfähigkeit von Arbeitssuchenden nicht unmittelbar steigern, ermöglichen sie es Arbeitgebenden, Rückschlüsse auf die Fähigkeiten des Arbeitssuchenden auf dem Arbeitsmarkt zu ziehen. Nach dem Signalmodell von Spence können Zertifikate als Signale für produktivitätsrelevante Eigenschaften wie Fähigkeiten, Motivation und Engagement fungieren und dadurch die Beschäftigungswahrscheinlichkeit von Arbeitssuchenden erhöhen \parencite[355--356]{Spence1973}.

\textcite{Kanas:2023} haben die Signaling-Theorie als einen möglichen theoretischen Ansatz zur Erklärung der Auswirkungen des Angebots an Sprachkursen auf die Integration in den Arbeitsmarkt herangezogen. Dieser Ansatz basiert auf der Annahme, dass Arbeitgebende über unvollständige Informationen über die Bewerbenden verfügen und daher auf Signale zurückgreifen. Auf Grundlage dieses theoretischen Ansatzes lassen sich, wie auch bei Kanas und Kosyakova, die Beschäftigungschancen von Geflüchteten in Deutschland erklären. In dem Artikel wird darauf hingewiesen, dass es für Arbeitgebende schwierig sein kann, die in den Herkunftsländern erworbenen Bildungsabschlüsse und ähnliche Zeugnisse von Geflüchteten zu bewerten. Daher stellen die im Aufnahmeland erworbenen Zertifikate bzw. Teilnahmebescheinigungen von Deutschkursen für Arbeitgebende besser beobachtbare und bewertbare Informationen dar. In Deutschland erworbene Sprachzertifikate liefern nicht nur ein verlässliches Signal hinsichtlich der Sprachkompetenzen von Geflüchteten, sondern ermöglichen es Arbeitgebenden auch, anhand dieser Zertifikate Rückschlüsse auf die Motivation und die Trainierbarkeit der Geflüchteten auf dem deutschen Arbeitsmarkt zu ziehen. Durch dieses Signal können Arbeitgebende eine positive Einschätzung des Arbeitssuchenden vornehmen, sodass dessen Beschäftigungschancen steigen könnten.

\section{Soziale Schließung und Credentialing}
\label{cha:Soziale Schließung}

„Credentialing“ ist ein theoretischer Ansatz, der von \textcite{collins1979credential} entwickelt wurde und auf der Annahme beruht, dass Zertifikate/Dokumente als Auslesemechanismus für eine bestimmte Arbeitsstelle dienen können. Nach \textcite{collins1979credential}’ Ansatz können Instrumente wie Ausbildung oder Sprachkurse zwar die Kenntnisse und Fähigkeiten des Individuums verbessern, im Arbeitsmarkt können Zertifikate/Dokumente jedoch eine Art Vorbedingung darstellen. Die Teilnahme bzw. Nichtteilnahme an einem Integrationskurs oder der Erwerb eines Zertifikats steht im Einklang mit Collins’ theoretischem Ansatz. Aus diesem Grund wird in dieser Arbeit der mögliche Einfluss der von Geflüchteten im Rahmen des Integrationskurses erworbenen Zeugnisse und Zertifikate auf die Beschäftigung anhand von Collins’ Credentialing-Theorie erläutert.

Die Credentialing-Theorie wird im Zusammenhang mit sozialer Schließung erläutert. Der Zugang zu bestimmten Stellen ist möglicherweise nicht für alle offen. Als Voraussetzung können Arbeitgebende bestimmte Ausbildungsnachweise, Sprachzertifikate und ähnliche Unterlagen verlangen. Diese Voraussetzung kann Personen, die nicht über die erforderlichen Zertifikate bzw. Unterlagen verfügen, von der Stellenbewerbung ausschließen; daher steht sie im Zusammenhang mit sozialer Schließung. Credentials bieten einer Gruppe Vorteile hinsichtlich der Mobilität, während die Mobilität der anderen Gruppe dadurch eingeschränkt wird \parencite[31--34]{collins1979credential}.



\section{Die Perspektive des Sozialkapitals}
\label{cha:sozialkapital}

\textcite{Kanas:2023} haben auf der Grundlage verschiedener Studien die Auswirkungen von Sprachkursen auf die Integration in den Arbeitsmarkt aus der Perspektive der Sozialkapitaltheorie erklärt. Das grundlegende theoretische Argument stützt sich darauf, dass durch Sprachkurse ein intensiverer Kontakt zwischen Flüchtlingen und der lokalen Bevölkerung hergestellt werden kann und dass dadurch die Integration der Geflüchteten in den Arbeitsmarkt verstärkt werden kann.

\textcite{doi:10.1111/imre.12130} haben versucht, die Frage zu beantworten, ob die Teilnahme an Sprachkursen durch türkische und marokkanische Zugewanderte langfristig Auswirkungen auf den Kontakt mit der lokalen Bevölkerung hat. Das zentrale Ergebnis der Studie zeigt, dass Sprachkurse nicht nur einen positiven Einfluss auf die Sprachkenntnisse haben, sondern auch einen positiven Zusammenhang mit dem Kontakt der Zugewanderten zur einheimischen Bevölkerung aufweisen. Aus diesem Grund können Sprachkurse Zugewanderten die Möglichkeit bieten, soziale Kontakte zur lokalen Bevölkerung aufzubauen.

\textcite{schuller2011integrationspanel} untersuchen die Ergebnisse der vom BAMF durchgeführten Integrationskurse. Die Studie befasst sich mit der Frage, inwieweit Integrationskurse für Zugewanderte nachhaltig und erfolgreich sind. In diesem Zusammenhang wurden Faktoren wie die Verbesserung der Deutschkenntnisse und die Häufigkeit der Kontakte zu Deutschen berücksichtigt und analysiert. \textcite{schuller2011integrationspanel} haben gezeigt, dass die Häufigkeit der Kontakte zu Deutschen zu Beginn des Kurses bei 3,4 lag und am Ende des Kurses bei 3,6 \parencite[233]{schuller2011integrationspanel}. Die Studie kam zu dem Ergebnis, dass die Teilnahme an Integrationskursen es zuwandernden Personen nicht nur ermöglicht, ihre Deutschkenntnisse zu verbessern, sondern auch die Häufigkeit der Kontakte zu Deutschen zu steigern.

\textcite{doi:10.1111/j.1747-7379.2010.00810.x} erläutern den Zusammenhang zwischen dem Kontakt zur lokalen Bevölkerung und der Integration in den Arbeitsmarkt. \textcite{doi:10.1111/j.1747-7379.2010.00810.x} haben festgestellt, dass der Kontakt von Geflüchteten in den Niederlanden zur lokalen Bevölkerung positiv mit deren Beschäftigungschancen zusammenhängt. Der Begriff „Bridging ties“ bezeichnet die sozialen Netzwerke, die Geflüchtete mit Menschen außerhalb ihrer eigenen Herkunftsgruppe aufbauen. Es wurde festgestellt, dass Geflüchtete mit einem höheren Anteil an „Bridging ties“ einen besseren beruflichen Status aufweisen.

Sprachkurse dienen nicht nur dem Spracherwerb, sondern können auch den Kontakt der Geflüchteten zur lokalen Bevölkerung stärken, wodurch diese einen besseren beruflichen Status erlangen können, wie dargelegt wurde. \textcite{Kanas:2023} haben im Rahmen dieser Untersuchungen versucht, zu erläutern, dass lokale Sprachkurse die Beschäftigungschancen von Geflüchteten durch die Schaffung von Kontaktmöglichkeiten zur lokalen Bevölkerung verbessern können. Da Sprachkurse den Geflüchteten die Möglichkeit bieten, die für den Kontakt mit der lokalen Bevölkerung erforderlichen Fähigkeiten zu erwerben, sind \textcite{Kanas:2023} zu der Vermutung gelangt, dass dies mit der Integration der Geflüchteten in den Arbeitsmarkt in Zusammenhang stehen könnte.


\section{Theoretische Erweiterung: Das handlungstheoretische Modell des Spracherwerbs (nach Esser)}
\label{cha:Erweiterung}

Die theoretische Erweiterung wird anhand des von \textcite{Esser:2006} entwickelten „Handlungstheoretischen Modells des Spracherwerbs“ erklärt. Das handlungstheoretische Modell des Spracherwerbs basiert auf der Vorstellung, dass das Erlernen einer Sprache eine Investition darstellt. Das Modell geht davon aus, dass Menschen für den Erwerb einer Sprache finanzielle Mittel, Zeit und persönlichen Aufwand einsetzen oder auf andere nutzenbringende Aktivitäten verzichten können. Esser weist darauf hin, dass das Erlernen einer Zweitsprache dem Individuum dadurch einen Nutzen bringen kann. Insbesondere bei Erwachsenen wird das Sprachenlernen als Investition betrachtet. Esser setzt den Spracherwerb mit bestimmten Aspekten des Humankapitals in Verbindung. Sprachkenntnisse stellen eine Eigenschaft oder Fähigkeit dar. Diese Fähigkeit lässt sich auf produktive Weise verwenden. Auf dem Arbeitsmarkt können diese Fähigkeiten beispielsweise produktiv genutzt werden. Außerdem erfordert der Spracherwerb Ressourcen. Aus diesem Grund beschreibt Esser das Erlernen einer Sprache als Investition, und darauf basiert der Ausgangspunkt seines Modells zum Spracherwerb \parencite[59--63]{Esser:2006}.

\textcite{Esser:2006} erklärt das Modell des Spracherwerbs anhand von vier verschiedenen Faktoren. Diese sind: Motivation, Zugang, Effizienz und Kosten. Der erste Baustein von Essers Investitionsmodell ist die Motivation. Die Motivation stellt den Unterschied zwischen dem Nutzen der neuen Sprache und dem aktuellen Nutzen der Muttersprache des Individuums dar. Im Modell wird der erwartete Nutzen des Erlernens der Zielsprache als U(L2) definiert, während der Nutzen der bestehenden Muttersprache des Individuums als U(L1) definiert wird. Die Motivation ist keine einzelne Zahl, sondern repräsentiert diesen Unterschied. Ein weiterer wichtiger Faktor für das Investitionsmodell sind die Strukturen des Zugangs bzw. der Zugangsmöglichkeiten. Der Kontakt des Menschen mit der Zielsprache im Alltag kann in diesem Zusammenhang beispielsweise als „Exposure“ gewertet werden. Auch Sprachkurse können in diesem Zusammenhang einen strukturierten Zugang zur Zielsprache ermöglichen. Ein weiterer Faktor ist die Effizienz. Diese wird definiert als der Anteil, in dem sprachlicher Input tatsächlich in Lernen umgewandelt wird. Als Determinanten der Effizienz werden beispielsweise das Einreisealter, die Intelligenz und die kulturelle Distanz angeführt. Schließlich fallen Kosten für den Lernprozess an. Zu den Kosten des Spracherwerbs zählen insbesondere Zeitkosten und andere Aufwendungen. Insbesondere beim gesteuerten Spracherwerb, etwa durch die Teilnahme an einem Sprachkurs, können dabei materielle Kosten und ein zusätzlicher zeitlicher Aufwand entstehen \parencite[82--89]{Esser:2006}. Das Verhältnis dieser vier Faktoren untereinander wird in der Gleichung~\eqref{eq:spracherwerb} des Investitionsmodells von \textcite{Esser:2006} dargestellt.

\begin{equation}
	U(L2) - U(L1) > C(L2) / p(exp) p(eff)
	\label{eq:spracherwerb}
\end{equation}

Die Differenz zwischen dem erwarteten Nutzen des Erwerbs einer Zweitsprache (L2) und dem Nutzen der bereits vorhandenen Sprache spiegelt die Motivation wider. C(L2) steht für die Kosten des Erwerbs der Sprache des Aufnahmelandes, p(exp) für die Möglichkeiten des Zugangs zur Zielsprache und p(eff) für die Effizienz des Spracherwerbs. Insgesamt erklärt Gleichung/Formel~\eqref{eq:spracherwerb}, dass bei geringerer Effizienz und eingeschränktem Zugang zum Spracherwerb ein höherer Motivationswert erforderlich ist.


\begin{figure}[htbp]
	\centering
	\resizebox{\textwidth}{!}{% <- keeps it inside the text block
		

	\begin{tikzpicture}[
	x=12cm, y=8cm,                    % <- overall size of the plot
	>={Latex[length=3mm,width=2mm]},
	font=\normalsize,
	case/.style={circle, fill, inner sep=2.4pt},
	tick/.style={circle, fill, inner sep=1.4pt}
	]
	
	%% ---------- adjustable parameters -------------------------------
	\def\Cm{0.22}      % cost constant of the low-cost curve  C(L2)_-
	\def\Cp{0.52}      % cost constant of the high-cost curve C(L2)_+
	\def\Mp{0.80}      % M_+  high motivation
	\def\Mz{0.44}      % M_0  medium motivation
	\def\Mm{0.14}      % M_-  low motivation
	\def\px{0.15}      % p    low  access x efficiency
	\def\pplus{0.82}   % p_+  high access x efficiency
	\def\ymax{1.30}    % top of the vertical axis
	
	%% ---------- axes -------------------------------------------------
	\draw[->] (0,0) -- (0,\ymax);
	\draw[->] (0,0) -- (1.16,0);
	
	\node[anchor=south west, align=left, inner sep=1pt] at (0.005,\ymax)
	{U(L2)-U(L1)\\ Motivation};
	\node[anchor=north west] at (1.085,-0.015) {$p(\mathrm{L2})$};
	\node[anchor=north]      at (0.45,-0.06)   {Zugang $\cdot$ Effizienz};
	
	%% ---------- tick marks -------------------------------------------
	\foreach \y/\lab in {\Mp/{$M_+$}, \Mz/{$M_0$}, \Mm/{$M_-$}}{
		\node[tick] at (0,\y) {};
		\node[anchor=east] at (-0.014,\y) {\lab};
	}
	\node[tick] at (\px,0)    {};
	\node[tick] at (\pplus,0) {};
	\node[anchor=north] at (0,-0.015)      {$0$};
	\node[anchor=north] at (\px,-0.015)    {$p$};
	\node[anchor=north] at (\pplus,-0.015) {$p_+$};
	\node[anchor=north] at (1,-0.015)      {$1$};
	
	%% ---------- the two cost curves  C/p ------------------------------
	% high-cost curve (dashed)
	\draw[dashed, dash pattern=on 8pt off 6pt, samples=400, smooth,
	domain={\Cp/\ymax}:1] plot (\x, {\Cp/\x});
	% low-cost curve (solid)
	\draw[samples=400, smooth,
	domain={\Cm/\ymax}:1] plot (\x, {\Cm/\x});
	
	\node at (0.50,{\Cp/0.50-0.13}) {$C_+/p$};
	\node at (0.50,{\Cm/0.50-0.09}) {$C_-/p$};
	
	%% ---------- right-hand end points, dotted vertical at p = 1 -------
	\draw[dotted, thick] (1,0) -- (1,\ymax);
	\node[tick] at (1,\Cp) {};
	\node[tick] at (1,\Cm) {};
	\node[anchor=west] at (1.02,\Cp) {$C(\mathrm{L2})_+$};
	\node[anchor=west] at (1.02,\Cm) {$C(\mathrm{L2})_-$};
	\node[anchor=west] at (1.02,{(\Cp+\Cm)/2}) {Kosten};
	
	%% ---------- region labels -----------------------------------------
	\node at (0.66,1.14) {Spracherwerb};
	\node[align=center] at (0.40,0.115) {kein\\ Spracherwerb};
	
	%% ---------- the six cases ------------------------------------------
	\foreach \x/\y/\n in {0.175/\Mp/1, 0.175/\Mz/2, 0.175/\Mm/3,
		\pplus/0.90/4, \pplus/\Mz/5, \pplus/\Mm/6}{
		\node[case] at (\x,\y) {};
		\node[anchor=east] at ({\x-0.022},\y) {\n};
	}
	
	%% ---------- outer frame ---------------------------------------------
	\draw ($(current bounding box.south west)+(-0.025,-0.02)$)
	rectangle
	($(current bounding box.north east)+(0.025,0.02)$);
	
\end{tikzpicture}


	

%
	}
	\caption[Motivation, Zugang, Effizienz und Kosten beim Spracherwerb]{%
		Das Zusammenspiel von Motivation, Zugang, Effizienz und Kosten beim
		(Zweit-)Spracherwerb. Eigene Darstellung in Anlehnung an
		% TODO: Seitenzahl bei Esser 2006 noch verifizieren
		\textcite[76]{Esser:2006}.}
	\label{fig:spracherwerb}
\end{figure}


Abbildung~\ref{fig:spracherwerb} veranschaulicht, wie die Entscheidung zum Sprachenlernen von Zugewanderten zustande kommt. Bei geringem Zugang zur Zielsprache und verminderter Effizienz wird ein höherer U(L2)-Bedarf erläutert, damit das Erlernen der Zweitsprache stattfinden kann. In ähnlicher Weise wird dargelegt, dass bei steigenden Kosten möglicherweise eine höhere Motivation erforderlich ist, um die Entscheidung zum Erlernen der Zielsprache zu treffen. Die Punkte in Abbildung~\ref{fig:spracherwerb} zeigen Situationen, in denen Spracherwerb je nach Kombination aus Motivation, Zugang und Effizienz möglich ist oder nicht. Außerdem erklären die verschiedenen Kurven die Kosten und deren Höhe. Demnach muss die Motivation für den Spracherwerb steigen, wenn die Kosten steigen.

\textcite{Esser:2006} verbindet die Motivation zum Erlernen einer Zweitsprache (L2) mit der Zeitperspektive sowie der Rückkehr- bzw. Bleibeperspektive. Im Zusammenhang mit der Zeitperspektive nach der Migration wurden Faktoren wie die Verwendbarkeit der Sprache und biografische Ziele als Faktoren genannt, die mit der Motivation zum Erlernen der Zielsprache in Zusammenhang stehen \parencite[82]{Esser:2006}. Es ist zu erwarten, dass die Motivation zum Erlernen einer Zweitsprache (L2) im Falle einer temporären Migration und einer Rückkehrorientierung abnehmen könnte.

Nach \textcite{Esser:2006} Modell des Spracherwerbs kann ein Sprachkurs den Zugang verbessern und die Effizienz steigern. Der Nutzen, den eine Person durch das Erlernen einer Sprache erzielt, kann jedoch auch mit der Motivation zusammenhängen. Aus diesem Grund ist zu erwarten, dass Flüchtlinge mit einem sicheren Aufenthaltsstatus hinsichtlich ihrer Bleibeperspektive tendenziell eine höhere Motivation aufweisen, in die Zielsprache (L2) zu investieren. Demgegenüber ist davon auszugehen, dass die Motivation zum Erlernen der Zielsprache bei Geflüchteten, deren Asylantrag abgelehnt wurde oder deren Asylantragsverfahren noch nicht abgeschlossen ist, tendenziell geringer ist. Insgesamt wird angenommen, dass sich die Geflüchteten je nach ihrem Aufenthaltsstatus unterscheiden lassen, welche Gruppe stärker vom lokalen Sprachkursangebot profitiert.

\textcite{Esser:2006} weist darauf hin, dass die Sprachkenntnisse von Geflüchteten auf dem Arbeitsmarkt verwendet werden können. Demnach kann davon ausgegangen werden, dass die Motivation zum Sprachenlernen bei Geflüchteten, deren Asylantrag anerkannt wurde, aufgrund ihrer Bleibeperspektiven tendenziell zunehmen könnte und dass sie stärker von lokalen Sprachkursen profitieren könnten. Infolgedessen wird davon ausgegangen, dass die erlernte Zielsprache auf dem Arbeitsmarkt produktiv verwendet werden kann \parencite[68--70]{Esser:2006}.
